% Copyright (c) 2018 by Daniel R. Driver.
% !TEX encoding = UTF-8 Unicode
% !TEX TS-program = XeLaTeX

\documentclass[titlepage]{article}

% This document presumes a file structure and set of inputs that are
% available at: git@github.com:danieldriver/syllabi.git

\newcommand\policy{../policy}
\newcommand\incl{../includes}
\ProvidesFile{variables.tex}[2022/09/11 v2.3 -- Syllabus variables]

\usepackage{xspace} % make manual spaces (like \mycmd\ ) unnecessary
\usepackage{xifthen} % provides \isempty test

% variables for internal use
\newcommand\prof{}
\newcommand\pdegree{}
\newcommand\pphone{}
\newcommand\pemail{}
\newcommand\poffice{}
\newcommand\phours{}
%
\newcommand\ccode{}
\newcommand\ctitle{}
\newcommand\cseries{}
\newcommand\cversion{}
\newcommand\csemester{}
\newcommand\cmeetson{}
\newcommand\cmeetsat{}
\newcommand\cmeetsin{}
\newcommand\cwebsite{}
\newcommand\cdescrip{}
\newcommand\cprereqs{}
\newcommand\edobject{}

% in case of fully online, BTh, or undergraduate courses - https://tex.stackexchange.com/a/5896
\newif\ifonline
\newcommand\Int[2]{\ifonline#1\else#2\fi}
\newif\ifbth
\newcommand\bth[2]{\ifbth#1\else#2\fi}
\newif\ifUG
\newcommand\UG[2]{\ifUG#1\else#2\fi}

% commands for setting variables in the preamble
\newcommand\professor[2][PhD]{
  \renewcommand\pdegree{#1\xspace}
  \renewcommand\prof{#2\xspace}}
\newcommand\phone[1]{
  \renewcommand\pphone{\addfontfeatures{Numbers=Monospaced}#1\xspace}}
\newcommand\email[1]{
  \renewcommand\pemail{\href{mailto:#1}{#1}\xspace}}
\newcommand\officehours[2][AST Library/Teams]{
  \renewcommand\poffice{#1\xspace}
  \renewcommand\phours{\href{https://outlook.office.com/bookwithme/user/b97b147a2e5c43b4a0e200c2cb630b77@astheology.ns.ca?anonymous}{#2}\xspace}%
  }
%
\newcommand\coursecode[2][1.0]{
  \renewcommand\cversion{#1\xspace}
  \renewcommand\ccode{#2\xspace}}
\newcommand\coursetitle[2][]{
  \ifthenelse{\isempty{#1}}%
    {}% do nothing if #1 is empty, else:
    {\renewcommand\cseries{#1\\[1ex]}}
  \renewcommand\ctitle{#2\xspace}}
\newcommand\semester[1]{
  \renewcommand\csemester{#1\xspace}}
\newcommand\meets[3]{
  \newcommand\AM{\textsc{am}}
  \newcommand\PM{\textsc{pm}}
  \renewcommand\cmeetson{#1\xspace}
  \renewcommand\cmeetsat{\Int{From 9:00 \AM}{#2}\xspace}
  \renewcommand\cmeetsin{\Int{\href{https://smu.brightspace.com/d2l/login}{Brightspace}}{#3}\xspace}}
\newcommand\website[2][]{
  \renewcommand\cwebsite{
  \ifthenelse{\isempty{#1}}%
    {\href{https://#2}}% copy url if #1 is empty, else
    {\href{#1}}% use specified url and print #2
  {#2}\xspace}}
\newcommand\cdescription[2][BF 1001; RM 1000 or GTRS 6000]{
  \renewcommand\cprereqs{#1}
  \renewcommand\cdescrip{#2\par}}
\newcommand\objectives[1]{
  \renewcommand\edobject{#1\par}}


%\onlinetrue % \Int{true}{false}
\coursecode[3.1]{HB 3113(Hyb)}
\coursetitle{The Ten Commandments}

% Taught as RLGS 3113 in:
%   - Winter 2012
% Taught as BSTH 3213 in:
%   - Winter 2015
% Taught as HB 3113(Int) in:
%   - Winter 2018
%   - Winter 2022

\professor{Daniel R. Driver}
\phone{902-425-7051}
\email{ddriver@astheology.ns.ca}
\officehours[Microsoft Teams]{Tuesdays, 1:00--3:00 \PM}

\semester{Winter Term 2022}
\meets{Tuesdays}% \meets{on}{at}{in}
      {4:00--6:30 \PM\ Atlantic}
      {Teams/Classroom 3}
\website{danieldriver.com}
\cdescription{% copy from the current Academic Calendar
	As the first words spoken by God at Sinai and the leading terms of a
	covenant sealed with blood, the Ten Commandments are well known as a
	premier instance of biblical law. This course examines their ancient
	origins. It also considers the weight they have held for Jews,
	Christians, and others through an exploration of the Decalogue’s
	reception history, in modes as diverse as liturgy, commentary, fine
	art, film, and popular culture.
%	From 2018 (truncated for Academic Calendar 2021--2022):
%
%	As the first words spoken by God at Sinai and the leading terms of a
%	covenant sealed with blood, the Ten Commandments are well known as
%	the preeminent instance of biblical law. Less well appreciated is
%	the complex legal and cultic context that gave rise to multiple
%	versions of the Decalogue in ancient Israel, as seen in Exodus 20,
%	Deuteronomy 5, Exodus 34, and Leviticus 19.
%
%	This course examines the ancient origins of the Ten Commandments as
%	a premier example of the Bible’s legal traditions. It also considers
%	the weight that these laws have held for Jews, Christians, and
%	others through an exploration of the Decalogue’s reception history,
%	which includes modes of expression as diverse as liturgy, biblical
%	commentary, fine art, film, and popular culture. Finally, with ample
%	time to consider a commandment per week, the course invites students
%	to reflect creatively on the ethical trajectory of the commandments,
%	their place in contemporary religious life, and their stature in and
%	beyond the Christian Church.
}% end of course description
\objectives{% recall Bloom's taxonomy: http://www.celt.iastate.edu/teaching/RevisedBlooms1.html
By the end of the course students should be able to:
	recite the Ten Commandments in a received order, and summarize the way different religious traditions count to Ten;
	locate and contextualize at least three versions of the Ten Commandments in the Bible;
	interpret each of the Commandments in their immediate literary contexts;
	integrate an understanding of the Ten Commandments with a broader view of biblical law;
	understand, in broad terms, how the Commandments have been interpreted in history;
	independently investigate the historic use of some Commandments;
	identify, evaluate, and report on contemporary connections with at least one of the Ten Commandments.
}% end of learning objectives

\ProvidesFile{preamble.tex}[2014/09/02 v1.1 -- Syllabus preamble]

% basic typography
\usepackage{fontspec,realscripts}
\defaultfontfeatures{Ligatures=TeX}
\setmainfont[%
  Numbers={Proportional,OldStyle},
  SmallCapsFeatures={RawFeature=+c2sc;+smcp,LetterSpace=1}]{Scala OT}
\newfontfamily\lining[Numbers={Lining}]{Scala OT}
\setsansfont[%
  SmallCapsFeatures={RawFeature=+c2sc;+smcp,LetterSpace=1}]{Scala Sans Pro}
\setmonofont[Scale=MatchLowercase]{Menlo}
\usepackage{sectsty}
\allsectionsfont{\sffamily\scshape}
\frenchspacing
\setlength{\emergencystretch}{3em} % prevent overfull lines

% custom font size and leading
\renewcommand\tiny{\fontsize{6}{9}\selectfont}
\renewcommand\scriptsize{\fontsize{7}{10}\selectfont}
\renewcommand\footnotesize{\fontsize{8}{10}\selectfont}
\renewcommand\small{\fontsize{9}{11}\selectfont}
\renewcommand\normalsize{\fontsize{10}{12}\selectfont}% base size
\renewcommand\large{\fontsize{11}{13}\selectfont}
\renewcommand\Large{\fontsize{12}{14}\selectfont}
\renewcommand\LARGE{\fontsize{16}{18}\selectfont}% "course syllabus \\ semester" benefits from more lead
\renewcommand\huge{\fontsize{20}{22}\selectfont}
\renewcommand\Huge{\fontsize{24}{26}\selectfont}

% layout packages: page, logo, tables
\usepackage[scale={0.6,0.8},
            xetex]{geometry}
\usepackage{graphicx}
\usepackage{array}     % allow insertions of column styling with >{}
\usepackage{booktabs}  % elegant horizontal rules in tables

% custom macros for a session count in the schedule of readings
\newcounter{session}
\newcounter{columns}
\newcounter{spancolumns}
\newcounter{midcolumns}
\newcounter{courseunit}
\newcommand\setcolumncount[2][0]{ % optionally set count to other than 0,
  \setcounter{session}{#1}        % e.g. to -1, or to a standing count
  \setcounter{columns}{#2}
  \setcounter{spancolumns}{#2-1}
  \setcounter{midcolumns}{#2-2}}
\newcommand\sessioncount{\stepcounter{session}\arabic{session}}
\newcommand\sessionskip[1]{\multicolumn{1}{@{}r@{ }}{#1}}
\newcommand\unit[1]{\multicolumn{\thecolumns}{c}{%
  \scshape\stepcounter{courseunit}\roman{courseunit}. #1}}
\newcommand\noclass[1]{\sessionskip{---.}&\multicolumn{\themidcolumns}{@{}l}{\itshape No Class: #1}}
\newcommand\rarr{\Large\char"EA31\hspace*{0.5pt}} % pointer finger w/ reminder tie in Scala OT
\newcommand\reminder[2]{\sessionskip{\rarr}&\multicolumn{\thespancolumns}{@{}l}{#1 \hfill #2}}

% color to match Tyndale's branding
\usepackage[usenames]{xcolor}
% predefined: black, white, red, green, blue, cyan, magenta, yellow
\definecolor{TyndaleURLs}{HTML}{0062A0} % links on tyndale.ca
\definecolor{TyndaleBlue}{cmyk}{1,1,0,.32}
\definecolor{TyndaleGold}{cmyk}{0,.27,1,0}
\definecolor{TyndaleRed}{cmyk}{0,1,.99,.04}
\definecolor{TyndaleBlack}{cmyk}{0,0,0,1}
\definecolor{TyndaleGreen}{cmyk}{.45,0,1,.24}
\definecolor{TyndaleOrange}{cmyk}{0,.79,1,0}
\definecolor{TyndaleAqua}{cmyk}{.47,0,.24,0}
\definecolor{TyndaleYellow}{cmyk}{.03,.03,.35,0}

% metadata (assumes a host of definitions are made in the main file)
\usepackage[setpagesize=false,     % leave this to geometry
            hyperfootnotes=false,  % fragile and distracting
            xetex]{hyperref}
\hypersetup{breaklinks=true,       % allow link text to break across lines
            colorlinks=true,       % colorlinks resets pdfborder to 0 0 0
            urlcolor=TyndaleURLs,  % for external links
            linkcolor=TyndaleRed,  % for normal internal links
            citecolor=TyndaleGold, % for bibliographical citations in text
            pdfauthor={\prof},
            pdftitle={\ccode: \ctitle},
            pdfsubject={Tyndale UC, \csemester},
            pdfcreator={github.com/danieldriver/syllabus}}
\urlstyle{same}                    % don't use monospace font for urls

% custom footlines
\usepackage{fancyhdr}
\pagestyle{fancy} % turn it on
\fancyhf{}        % reset everything
\renewcommand{\headrulewidth}{0pt} % remove header line as well
\lfoot{\sffamily\scshape\footnotesize\ctitle, \csemester}
\rfoot{\sffamily\scshape\footnotesize\prof\quad\thepage}

% language support: polyglossia needs to go last
\usepackage{polyglossia}
\setdefaultlanguage{english}
\setotherlanguage{hebrew}
\newfontfamily\hebrewfont[Script=Hebrew]{SBL BibLit}

% protect positioning of margin table in policy/grades -- needs to come after polyglossia
\usepackage{marginfix}

% gratuitous with custom title page, but useful as a fallback
\title{\ccode: \ctitle}
\author{\professor}
\date{\semester}


\begin{document}
\ProvidesFile{title.tex}[2013/09/06 v1.0 -- Syllabus title page]

\begin{titlepage}
  \begin{center}

    \begin{minipage}{\textwidth}
      \parbox[t]{0.5\textwidth}{
        \mbox{}\\[-\baselineskip] % dummy line to align parboxes
        \includegraphics[width=0.5\textwidth]{.syllabus/includes/TyndaleUC}}
      \hfill
      \parbox[t]{0.4\textwidth}{
        \raggedleft\LARGE\sffamily Course Syllabus\\
        \semester}
    \end{minipage}

    \vfill

    {\sffamily\textsc{\LARGE\MakeLowercase\coursecode}\\[2ex]
      \bfseries\itshape\Huge\coursetitle}

    \vfill

    \begin{tabular}{>{\bfseries}rl>{\bfseries}rl}
      \toprule
      Instructor & \professor                   & Course  & \liningnums{\coursecode}~\texttt{(v\version)} \\
      \midrule
      Phone      & \texttt{\phone}              & Meets   & \meeton                                       \\
      Email      & \href{mailto:\email}{\email} & Time    & \meetat                                       \\
      Office     & \officenumber                & Room    & \meetin                                       \\
      Hours      & \officehours                 & Website & \href{http://\website}{\website}              \\
      \bottomrule
    \end{tabular}

    \vfill

    \begin{description}
      \item[Commuter Hotline]
        Class cancellations due to inclement weather or illness will
        be announced on the commuter hotline at \texttt{416.226.6620
        x2187}. Alternately, weather cancellation information is posted
        at \href{http://tyndale.ca/weather}{tyndale.ca/weather}.
      \item[MyTyndale.ca]
        This course may have materials stored on its website, such as
        handouts or readings that may be needed in order to complete
        assignments. Students are responsible for checking these course
        pages on a regular basis. Here, too, students are able to view
        their grades throughout the semester. For more information see
        Section~\ref{mytyndale}, below.
      \item[Mail]
        Students are responsible for information communicated through
        their campus mailboxes and student e-mail accounts. A mailbox
        directory hangs beside the mailboxes. For more information
        contact the Registrar's office.
    \end{description}

  \end{center}

  \section{Course Description}
  \label{description}

  \emph{From the Academic Calendar:} \cdescription

\end{titlepage}
\setcounter{page}{2} % count the title page as page 1


  \section{Learning Objectives}
  \label{objectives}

  \edobject

\section{Required Texts \& Materials}
\label{texts}

The following texts are required. Students are strongly encouraged to
purchase their own copies. Visit the professor's website for links to
\href{https://danieldriver.com/courses/hb-3113/}{order the assigned editions}.

\begingroup
\renewcommand{\section}[2]{}% temporarily remove the section heading
\begin{thebibliography}{Centuries}% use the longest item in the bibliography

	\bibitem[Centuries]{Centuries} Jeffrey P. Greenman and Timothy Larsen, eds.
	\emph{The Decalogue through the Centuries: From the Hebrew Scriptures to Benedict XVI}.
	Louisville: Westminster John Knox, 2012.
	\textsc{isbn} 978-0664234904.

	\bibitem[Miller]{Miller} Patrick D. Miller.
    \emph{The Ten Commandments (Interpretation)}.
    Louisville: Westminster John Knox, 2009.
    \textsc{isbn} 978-0664230555.

	\bibitem[Rossetti]{Rossetti} Christina G. Rossetti.
	\emph{Letter and Spirit: Notes on the Commandments}.
	London: SPCK, n.d. [1883].
	Public domain: \url{https://archive.org/details/letterspiritnote00ross/}

\end{thebibliography}
\endgroup

If you do not have access to a good study Bible, I recommend either the
NRSV (Michael Coogan, ed.) or the NJPS (Adele Berlin and Marc Zvi
Brettler, eds.), both published by Oxford University Press. Both are also available through \href{http://ezproxy.astheology.ns.ca:2048/login?url=http://www.oxfordbiblicalstudies.com/}{Oxford Biblical Studies Online} (accessible with your AST student/faculty card barcode number and password).

\section{Supplementary Texts}
\label{supplementary}

The following titles may be referenced as supplementary texts. They are
not required. Other supplementary materials may be provided through the
course website.

\begingroup
\renewcommand{\section}[2]{}% temporarily remove the section heading
\begin{thebibliography}{3}% use the longest item in the bibliography

	\bibitem{BraatenSeitz} Carl E. Braaten and Christopher R. Seitz, eds.
	\emph{I Am the Lord Your God: Christian Reflections on the Ten Commandments}.
	Grand Rapids: Eerdmans, 2005.
	\textsc{isbn} 978-0802828125.

	\bibitem{Coogan} Michael Coogan.
    \emph{The Ten Commandments: A Short History of an Ancient Text}.
    New Haven: Yale University Press, 2014.
    \textsc{isbn} 978-0300178715.

	\bibitem{VanHarn} Roger E. Van Harn, ed.
	\emph{The Ten Commandments for Jews, Christians, and Others}.
	Grand Rapids: Eerdmans, 2007.
	\textsc{isbn} 978-0802829658.

	\bibitem[Making Sense]{rlgs} Northey, Margot, Bradford A. Anderson, and Joel N. Lohr.
	\emph{Making Sense in Religious Studies: A Student's Guide to Research and Writing}.
	3rd ed. Toronto: Oxford University Press, 2019. ISBN 978-0199026838.

	\bibitem[SBL2]{sbl2} Collins, Billie Jean, et al.
	\emph{The SBL Handbook of Style}.
	2nd ed. Atlanta: SBL Press, 2014. ISBN 978-1589839649.
	Designed to augment \href{http://www.chicagomanualofstyle.org/home.html}{\emph{Chicago Style}}
	(the standard at AST), there is also a free
	\href{https://www.sbl-site.org/assets/pdfs/pubs/SBLHSsupp2015-02.pdf}{Student Supplement for SBL2}.


\end{thebibliography}
\endgroup

\section{Course Outline}
\label{outline}

We will adhere to the schedule in \autoref{schedule} as closely as
possible, though the professor reserves the right to adjust it to suit
the needs of the class.

Note that readings are listed by the textbook's key word and chapter
numbers, except in the case of Rossetti, which gives page numbers.

\setcolumncount{6}% set up \sessioncount, \unit{}, \noclass{}, and \reminder{memo}{date} macros
\begin{table}[htbp]% add p to put the schedule on its own page
  \centering
  \begin{tabular}{>{\sessioncount.}r@{ }llllr}% make sure the column config agrees with \setcolumncount
	\toprule
	\sessionskip{\textbf{\S.}}&\textbf{Command}&\textbf{Exposition}&\textbf{Reception}&\textbf{Rossetti}&\textbf{Date}\\
	\midrule
		& (Intro)  &              & Centuries 11    &             & 11 Jan. \\%
		& First    & Miller Intro & Centuries 1     & pp. 6–25    & 18 Jan. \\% https://archive.org/details/letterspiritnote00ross/page/n7/mode/2up
		& Second   & Miller 1     & Centuries 2–3   & pp. 26–45   & 25 Jan. \\%
		& Third    & Miller 2     & Centuries 4     & pp. 46–65   &  1 Feb. \\%
		& Fourth   & Miller 3     & Centuries 5     & pp. 66–85   &  8 Feb. \\%
		& Fifth    & Miller 4     & Centuries 6     & pp. 86–105  & 15 Feb. \\%
	\noclass{Reading Week}                                        & 22 Feb. \\%
		& Sixth    & Miller 5     & Centuries 7     & pp. 106–125 &  1 Mar. \\%
		& Seventh  & Miller 6     & Centuries 8     & pp. 126–145 &  8 Mar. \\%
		& Eighth   & Miller 7     & Centuries 9     & pp. 146–165 & 15 Mar. \\%
		& Ninth    & Miller 8     & Centuries 10    & pp. 166–185 & 22 Mar. \\
		& Tenth    & Miller 9     & Centuries 12–13 & pp. 186–206 & 29 Mar. \\
		& (Recap)  & Miller end   & Film: tbd       &             & 5 Apr. \\
	\bottomrule
  \end{tabular}
  \caption{Schedule of Readings}
  \label{schedule}
\end{table}

See the AST website for a list of other \href{http://www.astheology.ns.ca/students/academic-dates.html}{important dates}.

\section{Evaluation}
\label{evaluation}

Given the small class size this term, I have adopted a tutorial format
for \ccode.

\begin{enumerate}

	\item Each week you will write a 1000–1200 word essay responding to
	a prompt.

	\item Essays should interact thoughtfully with both primary and
	secondary sources, and should be \textbf{thesis-driven}. Your essay
	must argue a thesis, and not simply describe what you have learned
	or what others have argued. At the tutorial, you should know what
	you think about the question, and should be able to summarize the
	argument of your essay in one sentence.

	\item Reading preparation should always give priority of attention
	to the assigned primary texts. Secondary literature will be
	important in lending guidance to the essay, but scholarly opinions
	should always be assessed on the basis of how helpful they are in
	explaining the primary text.

	\item By email or Teams, please send me your essay as a Microsoft
	Word document by 4pm the day before class. Essays received on time
	will receive written feedback. Late essays are not guaranteed
	written feedback. If you do not provide an essay at all, class may
	be cancelled. Essays will not be accepted after class (with rare
	exceptions).

	\item Written feedback will be returned by email or on Teams
	following discussion in the tutorial. Marks will not be assigned to
	the essays initially, and only in two cases later. Instead,
	affirmation and redirection will be provided as needed.

	\item Select two of your best essays to revise and resubmit for a
	formal mark. Both essays must be submitted before the last day of
	class. You can opt to submit them early.

%	Assignments from 2018:
%
%	\item I will set \textbf{discussion prompts} for each of the Ten
%	Commandments. They are intended to help you engage with the assigned
%	reading material, and with your classmates. Prompt responses should
%	be about 250 words long. Questions will go up by the start of each
%	class (10 \AM\ Atlantic Time). Answer the week's question by the end
%	of Tuesday. Then, by the end of Wednesday, read through the
%	responses of your classmates, and post another 50 words or so as you
%	interact with them.
%
%	\item Students are to make a \textbf{creative presentation} on one
%	of the Ten Commandments in the week in which that Commandment is
%	under investigation. Week 3 is an exception, since it is reserved
%	for the Prologue and Commandments I and II. Presentations may take
%	the form of a podcast or video, and their contents may be anything
%	that merits the adjective “creative” -- an interview, a dramatic
%	piece, some form of engagement with culture, or anything else
%	suitable that you envision. Each one should last 15 to 20 minutes,
%	and should feel like it deserves 20\% of the final grade. The only
%	strict requirement is that the work be sharable by URL. (Try
%	SoundCloud or YouTube.)
%
%	\item A \textbf{review essay} on Coogan's \emph{The Ten
%	Commandments} will be due at the end of the sixth week of class
%	(Friday, 16 February 2018). It should be 2,250 words long, plus or
%	minus 10\%. The first third of the essay should be a summary of the
%	work, while the remaining two thirds should be devoted to critical
%	engagement and evaluation.
%
%	\item A comprehensive \textbf{final paper} will give students an
%	opportunity to explore one of the Ten Commandments in detail,
%	complete with a thesis that involves original exposition and
%	theological reflection. It is due the Friday before the last class
%	(6 April 2018). It should be 3,500 words long, plus or minus 10\%
%	(the total count includes footnotes but not the final bibliography).

\end{enumerate}

The breakdown for the semester's total work is shown in
\autoref{grade-dist}.

\begin{table}[htbp]
  \centering
  {\lining
  \begin{tabular}{lr}
    \toprule
    Completion of All Essays      & 70\% \\
    Assessment of Selected Essays & 30\% \\
    \bottomrule
  \end{tabular}}
  \caption{Distribution of Grades}
  \label{grade-dist}
\end{table}

\ProvidesFile{grades.tex}[2023/09/06 v2.5 -- Course policy]

AST's \href{https://www.astheology.ns.ca/students/resources.html}{Academic
Calendar} provides guidelines and criteria for academic assessment. Marks are
assigned by letter grade using these benchmarks. Note that graduate and
undergraduate programs are currently assessed on a different scales at AST (see
\autoref{grade-syst}). For graduate students, a grade of C is deemed
``Unsatisfactory,'' whereas for undergraduate (BTh) students, a C is
``Satisfactory'' and the grade of D is possible for ``Marginal'' performance.
Other grades have different definitions, too.

\begin{table}[htbp]
 \centering
 {\lining\footnotesize
 \begin{tabular}{lll|ll}
   \toprule
               & \multicolumn{2}{c|}{Graduate Courses}   & \multicolumn{2}{c}{Undergraduate Courses} \\
               \cline{2-5}
   A+          & 94--100    & Exceptional               & 90--100    & Excellent                 \\
   A           & 87--93     & Outstanding               & 85--89     & Excellent                 \\
   A\char"2212 & 80--86     & Excellent                 & 80--84     & Excellent                 \\ [1ex]
   B+          & 77--79     & Good                      & 75--79     & Good                      \\
   B           & 73--76     & Acceptable                & 70--74     & Good                      \\
   B\char"2212 & 70--72     & Marginally Acceptable     &            &                           \\ [1ex]
   C           & 60--69     & Unsatisfactory            & 60--69     & Satisfactory              \\
   D           &            &                           & 50--59     & Marginal                  \\
   F           & 0--59      & Failure                   & 0--49      & Failure                   \\
   FP          & 0          & Failure due to Plagiarism & 0          & Failure due to Plagiarism \\
   \bottomrule
 \end{tabular}}
 \caption{Outline of Grades at AST}
 \label{grade-syst}
\end{table}

\ProvidesFile{other.tex}[2018/0/03 v2.4 -- Course policy]

\section{Other Course Policy}
\label{policy}

Late work will not be accepted, except in genuinely extenuating
circumstances. Students must submit something before the deadline if
they wish to receive credit. Unless I state otherwise, written
assignments are to be uploaded by 11:59 \PM\ on the date indicated.

Essay submissions must be typewritten, double-spaced, and formatted as
PDFs. They should be free from error. In this course they should follow
either \emph{The SBL Handbook of Style} (refer to the online
\href{https://www.sbl-site.org/assets/pdfs/pubs/SBLHSsupp2015-02.pdf}{
Student Supplement}), or \emph{The Chicago Manual of Style}, per AST's
policy for major papers. As a reminder, AST also upholds an Inclusive
Language Policy.

\href{http://www.eerdmans.com/Pages/Item/59043/Commentary-Statement.aspx}{Plagiarism},
if \href{https://www.theguardian.com/world/2013/feb/09/german-education-minister-quits-phd-plagiarism}{detected},
will result in failure of the course.

Students should request permission to record a class or lecture. If
permission is granted, or if recordings are provided (as in the case of
an online course), I stipulate that all recordings be for personal use
only. They may not be shared or distributed.

If you have abilities or disabilities that require modifications to the
assessment process or other aspects of this course, please advise the
course instructor as soon as possible.

Finally, I encourage the conscientious use of laptops, tablets, and
other technology in my classes. In classroom settings, realize that, as
\href{http://dx.doi.org/10.1016/j.compedu.2012.10.003}{cognitive
psychologists have demonstrated}, ``laptop multitasking hinders
classroom learning for both users and nearby peers.'' Do your part to
foster an environment of open dialogue by honouring the presence of your
classmates. In online settings, consider both the physical environment
in which you choose to work and the virtual environment that you help
create through your participation in various forums. Let your engagement
in this course be marked by rigour and charity alike.

\end{document}
